\documentclass[11pt]{letter}
\usepackage[utf8]{inputenc}
\usepackage{fontspec}
\usepackage{geometry}
\geometry{a4paper, margin=2.5cm}
\usepackage{hyperref}
\hypersetup{colorlinks, urlcolor=blue}
\usepackage{parskip}

\address{
  Rodrigo Tertulino \\
  Coordenação de Informática, Campus Mossoró \\
  Instituto Federal de Educação, Ciência e Tecnologia do Rio Grande do Norte (IFRN) \\
  Av. Presidente Dutra, s/n, Alto de São Manoel \\
  Mossoró, RN, 59628-330, Brazil \\
  \href{mailto:rodrigo.tertulino@ifrn.edu.br}{rodrigo.tertulino@ifrn.edu.br}
}

\signature{Rodrigo Tertulino \\ (on behalf of all authors)}

\begin{document}

\begin{letter}{The Editors-in-Chief \\
Journal of Internet Services and Applications (JISA) \\
Brazilian Computing Society (SBC)}

\opening{Dear Editors,}

We are pleased to submit our original research article entitled \textbf{``A Decentralized Identity Protocol for Trustworthy Federated Learning: Comparative Analysis of Trust Models for Internet-Based Health Services''} for consideration for publication in the \textit{Journal of Internet Services and Applications (JISA)}.

\section*{Fit with JISA Scope}

This work directly addresses JISA's core focus areas of \textbf{security and privacy}, \textbf{networking protocols and architectures}, and \textbf{applications} of Internet-based services. We design and empirically validate a decentralized identity protocol based on W3C standards (DIDs, VCs) that secures Federated Learning systems against Sybil attacks in Internet-distributed environments. The protocol operates as an Internet middleware layer---leveraging blockchain smart contracts for decentralized authorization and IPFS for content-addressed model storage---providing deterministic security guarantees without reliance on centralized trust anchors. The clinical healthcare setting serves as the motivating application domain, representative of any Internet-based service handling sensitive, regulated data.

\section*{Main Contributions}

We believe this manuscript makes the following original contributions to the literature on Internet services and applications:

\begin{enumerate}
  \item \textbf{A formal four-level taxonomy of trust models} (T0--T3) for Federated Learning, characterizing each model along verification mechanism, security guarantee type, LGPD regulatory alignment, and operational cost---providing a systematic framework absent from prior work.

  \item \textbf{The SSI-FL protocol:} a decentralized participant authentication protocol integrating W3C Decentralized Identifiers (DIDs), Verifiable Credentials (VCs), and blockchain smart contracts, specified in a domain-agnostic form applicable to any Internet-based collaborative machine learning scenario beyond healthcare.

  \item \textbf{Empirical validation on a clinical dataset (MIMIC-IV)} demonstrating that the T3 trust model neutralizes 100\% of external Sybil attacks, preserves clinically meaningful predictive performance (AUC-ROC = 0.954, Recall = 0.890), and incurs protocol overhead below 1\% of training time.

  \item \textbf{An analytical cost model} $C(N,R)$, calibrated from experimental data on the Ethereum Sepolia testnet, enabling prospective consortia to project operational costs before deployment (e.g., \textasciitilde\$15.67 for $N=10$ participants over $R=100$ rounds).

  \item \textbf{A structured LGPD compliance mapping} across all four trust models, the first such analysis in the BFL literature, providing actionable guidance for Brazilian institutions operating Internet-based health data processing services.
\end{enumerate}

\section*{Novelty and Significance}

To our knowledge, this is the first work to combine proactive W3C SSI-based identity verification with Federated Learning, and the first BFL study to include an analytical cost model, attack tree analysis, and LGPD compliance mapping within a single evaluation framework. The systematic literature review by Qammar et al.\ (2023) explicitly identified participant authentication as an unresolved open problem in BFL---our work responds directly to this gap.

The protocol's negligible overhead ($< 1\%$ of training time) and linear cost scaling with consortium size remove the two most commonly cited barriers to blockchain-integrated FL adoption in production Internet services.

\section*{Manuscript Information}

\begin{itemize}
  \item \textbf{Title:} A Decentralized Identity Protocol for Trustworthy Federated Learning: Comparative Analysis of Trust Models for Internet-Based Health Services
  \item \textbf{Authors:} Rodrigo Tertulino, Ricardo Almeida, Laercio Alencar (all affiliated with IFRN, Brazil)
  \item \textbf{Keywords:} Federated Learning, Trust Models, Self-Sovereign Identity, Decentralized Identifiers, Verifiable Credentials, Blockchain, LGPD, Internet Security, Privacy, Sybil Attack
  \item \textbf{Article type:} Research article
  \item \textbf{Number of figures:} 3 (color conceptual figures)
  \item \textbf{Number of tables:} 6
\end{itemize}

\section*{Declarations}

\begin{itemize}
  \item This manuscript has not been previously published and is not under consideration by any other journal.
  \item All authors have read and approved the final version of the manuscript.
  \item The authors declare no competing interests.
  \item The source code is publicly available at \url{https://github.com/rodrigoronner/TBFL-EHR-Framework}.
  \item The MIMIC-IV dataset was accessed under a PhysioNet credentialed user agreement.
  \item AI-assisted tools were used for language editing; the authors take full responsibility for the scientific content.
\end{itemize}

We are confident that this manuscript aligns closely with JISA's mission to publish recent advances in Internet-related science and technology, particularly in security, protocols, and applications. We would be honored to have it considered for publication.

\closing{Respectfully submitted,}

\end{letter}

\end{document}
